\documentclass[12pt]{article}

\usepackage[T1]{fontenc}
\usepackage[utf8]{inputenc}
\usepackage[brazil]{babel}
\usepackage{lmodern}
\usepackage{microtype}
\usepackage[a4paper,margin=2.5cm]{geometry}
\usepackage{setspace}
\usepackage{indentfirst}

\setstretch{1.15}

\title{ASPECTOS MAIS AMPLOS DOS CICLOS ECONÔMICOS \thanks{Texto Traduzido livremente por Luiz Mario Andrade com o auxílio da ferramenta Chat GPT 5.4}}
\author{Wesley C. Mitchell}
\date{1913}

\begin{document}

\maketitle

CAPÍTULO 5

ASPECTOS MAIS AMPLOS DOS CICLOS ECONÔMICOS

I. Resumo da Teoria Precedente dos Ciclos Econômicos

A TEORIA dos ciclos econômicos apresentada nos Capítulos 1-4 é uma análise descritiva dos processos de mudança cumulativa pelos quais uma retomada da atividade se transforma em prosperidade intensa, pela qual essa prosperidade gera uma crise, pela qual a crise se transforma em depressão e pela qual a depressão, após tornar-se mais severa por um tempo, finalmente leva a uma retomada da atividade semelhante àquela com que o ciclo começou.

Esta análise repousa primordialmente sobre uma elaborada investigação estatística dos fenômenos de ciclos recentes nos Estados Unidos, Grã-Bretanha, França e Alemanha. A linha de ataque estatística foi escolhida porque o problema é essencialmente quantitativo, envolvendo a importância relativa de diversas forças que são, elas mesmas, os resultados líquidos de inumeráveis decisões de negócios. A seleção dos dados estatísticos, os métodos de apresentação e a coordenação dos resultados foram determinados em grande parte por ideias emprestadas de autores teóricos ou de periódicos financeiros. Mas todas as tabelas e todas as ideias emprestadas foram ajustadas em uma estrutura fornecida por um estudo da organização econômica de hoje, que mostrou que o processo industrial de fabricação e o processo comercial de distribuição de bens são ambos inteiramente subordinados ao processo de negócios de ganhar dinheiro.

A teoria derivada desses materiais preencheu tantas páginas que um resumo pode ser útil. Há sempre o perigo, entretanto, de que um resumo plausível receba peso excessivo. Leitores familiarizados com teorias de ciclos econômicos apreciarão como é fácil fazer com que muitas explicações diferentes de crises pareçam convincentes quando a atenção se limita a uma gama restrita de fenômenos. Somente submetendo qualquer teoria ao teste prático de explicar a experiência real dos negócios é que seu valor pode ser determinado. O argumento a favor da presente teoria, portanto, e também o argumento contra ela, deve ser encontrado não no resumo fácil que se segue, mas nos capítulos difíceis que precedem, incluindo as Partes I e II, ou, melhor ainda, em um esforço independente de usá-la na interpretação do fluxo e refluxo incessante da atividade econômica.

1. CUMULAÇÃO DA PROSPERIDADE

Com qualquer fase do ciclo econômico que a análise comece, ela deve tomar como pressupostas as condições trazidas pela fase anterior, adiando a explicação dessas suposições até que tenha percorrido todo o ciclo e retornado ao seu ponto de partida.

Uma retomada da atividade, então, começa com este legado da depressão: um nível de preços baixo em comparação com os preços da prosperidade, reduções drásticas nos custos de fazer negócios, margens de lucro estreitas, reservas bancárias liberais, uma política conservadora na capitalização de empresas comerciais e na concessão de créditos, estoques moderados de mercadorias e compras cautelosas.

Por razões que aparecerão na sequência, tais condições são acompanhadas por uma expansão no volume físico do comércio. Embora lenta no início, esta expansão é cumulativa. Agora é apenas uma questão de tempo para que um aumento no montante de negócios transacionados, que cresce mais rapidamente à medida que avança, transforme o marasmo em atividade. Deixada por conta própria, esta transformação é efetuada por graus lentos; mas é frequentemente acelerada por algum evento propício oriundo de fontes que não os negócios domésticos, tais como colheitas excepcionalmente lucrativas, compras pesadas de suprimentos pelo governo ou um aumento acentuado na demanda de exportação para os produtos da indústria nacional.

Mesmo quando uma retomada da atividade está confinada inicialmente a uma gama estreita de indústrias ou a uma única seção do país, ela logo se espalha para outras partes do campo de negócios. Pois as empresas ativas devem comprar mais materiais, mercadorias e suprimentos correntes de outras empresas, estas últimas de outras ainda, e assim por diante, sem limites atribuíveis. Enquanto isso, todas as empresas que se tornam mais ocupadas empregam mais mão de obra, usam mais dinheiro emprestado e obtêm lucros maiores. As rendas familiares aumentam e a demanda dos consumidores expande-se, da mesma forma, em círculos cada vez mais amplos. Os lojistas repassam pedidos maiores de bens de consumo para comerciantes atacadistas, fabricantes, importadores e produtores de matérias-primas. Todas essas empresas exigem mais suprimentos de vários tipos para lidar com seu comércio crescente e aumentam as somas que pagam a funcionários, credores e proprietários — estimulando assim novamente a demanda por bens de produção e de consumo. Cedo ou tarde, essa expansão de pedidos chega às empresas das quais o ímpeto para maior atividade foi recebido pela primeira vez, e então toda essa série complicada de reações recomeça em um nível mais alto de intensidade. Durante todo esse tempo, a retomada da atividade está instilando um sentimento de otimismo entre os empresários, e esse sentimento se justifica e aumenta as forças que o geraram, tornando todos mais prontos para comprar livremente.

Embora o nível de preços esteja frequentemente caindo lentamente quando uma retomada começa, a expansão cumulativa no volume físico do comércio acaba por interromper a queda e inicia uma alta. Pois, quando as empresas têm em vista tantos negócios quanto podem suportar com suas instalações existentes de eficiência padrão, elas exigem preços mais altos para pedidos adicionais. Essa política prevalece mesmo nos ramos mais competitivos, porque pedidos adicionais só podem ser executados mediante a contratação de novos trabalhadores, o acionamento de máquinas antigas, a compra de novos equipamentos ou alguma outra mudança que acarrete aumento de despesas. A expectativa de sua vinda acelera o avanço. Os compradores estão ansiosos para obter ou contratar grandes suprimentos enquanto o baixo nível de cotações continua, e os primeiros sinais definitivos de uma tendência ascendente de cotações trazem uma corrida repentina de pedidos.

Assim como o aumento no volume físico dos negócios, a alta nos preços se espalha rapidamente; pois cada avanço nas cotações pressiona alguém a se ressarcir fazendo um avanço compensatório nos preços do que tem a vender. As mudanças resultantes nos preços estão longe de ser uniformes, não apenas entre diferentes mercadorias, mas também entre diferentes partes do sistema de preços. Os preços de varejo ficam atrás dos preços de atacado, os preços de bens de consumo básicos atrás dos preços de bens de produção básicos, e os preços de produtos acabados atrás dos preços de suas matérias-primas. Entre as matérias-primas, os preços dos produtos minerais refletem as mudanças nas condições de negócios mais regularmente do que os preços dos produtos animais, agrícolas ou florestais. Os salários sobem frequentemente de forma mais imediata, mas sempre menos do que os preços de atacado; as taxas de desconto sobem às vezes mais lentamente do que as mercadorias e às vezes mais rapidamente; as taxas de juros em empréstimos de longo prazo movem-se sempre vagarosamente nos estágios iniciais da retomada, enquanto os preços das ações — particularmente das ações ordinárias — precedem e excedem os preços das mercadorias na subida.

As causas dessas diferenças na prontidão e na energia com que várias classes de preços respondem ao estímulo da atividade econômica encontram-se em parte em diferenças de organização nos mercados de mercadorias, trabalho, empréstimos e títulos; em parte nas circunstâncias técnicas que afetam a demanda e oferta relativas dessas diversas classes de bens; e em parte no ajuste dos preços de venda às mudanças no conjunto de preços de compra que uma empresa paga, em vez de mudanças nos preços dos bens específicos comprados para revenda.

Na grande maioria das empresas, lucros maiores resultam dessas flutuações divergentes de preços combinadas com vendas maiores. Pois, enquanto os preços das matérias-primas e das mercadorias compradas para revenda geralmente, e os preços dos empréstimos bancários frequentemente, sobem mais rápido do que os preços de venda, os preços da mão de obra ficam muito atrás, e os preços que compõem os custos suplementares são principalmente estereotipados por um tempo por antigos contratos relativos a salários, aluguéis e títulos.

Este aumento nos lucros, combinado com a prevalência do otimismo nos negócios, leva a uma marcada expansão dos investimentos. Naturalmente, as pesadas encomendas de máquinas, os grandes contratos para novas construções e assim por diante, que resultam disso, inflam ainda mais o volume físico dos negócios e tornam ainda mais fortes as forças que estão impulsionando os preços para cima.

De fato, a característica saliente desta fase do ciclo econômico é o funcionamento cumulativo dos vários processos que estão convertendo uma retomada do comércio em prosperidade intensa. Não apenas cada aumento no comércio causa outros aumentos, cada convertido ao otimismo faz novos convertidos, e cada avanço nos preços fornece um incentivo para novos avanços; mas o crescimento do comércio também ajuda a espalhar o otimismo e a elevar os preços, enquanto o otimismo e a alta dos preços se apoiam mutuamente e estimulam o crescimento do comércio. Finalmente, como acaba de ser dito, as mudanças em andamento nesses três fatores inflam os lucros e incentivam os investimentos, enquanto lucros altos e investimentos pesados reagem aumentando o comércio, justificando o otimismo e elevando os preços.

2. COMO A PROSPERIDADE GERA A CRISE

Enquanto os processos esboçados trabalham cumulativamente por um tempo para aumentar a prosperidade, eles também causam uma lenta acumulação de tensões dentro do sistema equilibrado de negócios — tensões que, em última análise, minam as condições sobre as quais a prosperidade repousa. Entre elas está o aumento gradual nos custos de fazer negócios. O declínio nos custos suplementares por unidade de produção cessa quando as empresas já reservaram todos os negócios que podem suportar com seu equipamento padrão, e um lento aumento nesses custos começa quando a expiração de contratos antigos força renovações às altas taxas de juros, aluguel e salários que prevalecem na prosperidade. Enquanto isso, os custos primários sobem a uma taxa relativamente rápida. Equipamentos e plantas antiquadas, mal localizadas ou que de outra forma trabalham com desvantagem, são postos novamente em operação. O preço da mão de obra sobe, não apenas porque as taxas padrão de salários aumentam, mas também porque o pagamento por horas extras é maior. Mais sério ainda é o declínio na eficiência da mão de obra, porque as horas extras trazem cansaço, devido ao emprego de "indesejáveis" e porque as equipes não podem ser conduzidas à velocidade máxima quando os empregos são mais numerosos do que os homens para preenchê-los. Os preços das matérias-primas continuam a subir mais rápido, em média, do que os preços de venda dos produtos. Finalmente, os numerosos pequenos desperdícios, incidentes à condução das empresas comerciais, aumentam quando os gerentes estão apressados por uma pressão de pedidos exigindo entrega imediata.

Uma segunda tensão é a acumulação de tensão nos mercados de investimento e monetário. A oferta de fundos disponíveis às taxas de juros anteriores para a compra de títulos, para empréstimos hipotecários e afins, deixa de acompanhar a demanda em rápido crescimento. Torna-se difícil negociar novas emissões de títulos, exceto em termos onerosos, e os homens de negócios queixam-se da "escassez de capital". Nem a oferta de empréstimos bancários cresce rápido o suficiente para acompanhar a demanda. Pois a oferta é limitada pelas reservas que os banqueiros mantêm contra seus passivos de demanda em expansão. O pleno emprego e o comércio varejista ativo fazem com que tanto dinheiro permaneça suspenso na circulação ativa que o dinheiro deixado nos bancos aumenta de forma bastante lenta, mesmo quando a produção de ouro está subindo mais rapidamente. Por outro lado, a demanda por empréstimos bancários cresce não apenas com o volume físico do comércio, mas também com a alta dos preços e com o desejo dos homens de negócios de usar seus próprios fundos para controlar o maior número possível de empreendimentos. Além disso, essa demanda é relativamente inelástica, uma vez que muitos tomadores de empréstimos pensam que podem pagar altas taxas de desconto por alguns meses e ainda obter lucros sobre seu faturamento, e uma vez que as corporações que não estão dispostas a vender títulos de longo prazo nos termos severos que passaram a prevalecer tentam levantar parte dos fundos de que necessitam descontando notas de um ou dois anos.

A tensão nos mercados de títulos e monetário é desfavorável à continuação da prosperidade, não apenas porque as altas taxas de juros reduzem as margens de lucro prospectivas, mas também porque elas freiam a expansão no comércio a partir da qual a prosperidade se desenvolveu. Muitos empreendimentos projetados são abandonados ou adiados, seja porque os tomadores de empréstimo concluem que os juros absorveriam grande parte de seus lucros, seja porque os credores se recusam a estender ainda mais seus compromissos.

Um grupo importante de empresas sofre um revés especialmente severo por esta causa em conjunção com preços altos — o grupo que depende primordialmente da demanda por equipamentos industriais. Nos estágios iniciais da prosperidade, este grupo geralmente desfruta de uma temporada de atividade excepcionalmente intensa. Mas quando o mercado de títulos torna-se restrito e — o que é frequentemente mais importante — quando o custo de construção tornou-se alto, empresas comerciais e capitalistas individuais adiam a execução de muitos planos para estender plantas antigas e erguer novas. Como resultado, os contratos para este tipo de trabalho tornam-se menos numerosos à medida que o clímax da prosperidade se aproxima. Então as siderúrgicas, fundições, fábricas de máquinas, fundições de cobre, pedreiras, serrarias, fábricas de cimento, empresas de construção, empreiteiras gerais e similares veem seus pedidos para entrega futura diminuírem. Enquanto por enquanto podem estar trabalhando em alta pressão para completar contratos antigos dentro do período estipulado, elas enfrentam uma séria restrição de comércio no futuro próximo.

O imponente edifício da prosperidade é construído com um generoso fator de segurança; mas quanto maior cresce a estrutura, mais severas tornam-se essas tensões internas. O único meio eficaz de prevenir o desastre enquanto se continua a construir é elevar os preços de venda vez após vez o suficiente para compensar as invasões dos custos sobre os lucros, para anular as taxas de juros crescentes e para manter os investidores dispostos a contratar novos equipamentos industriais.

Mas é impossível manter os preços de venda subindo indefinidamente. Na falta de outros freios, a inadequação das reservas de caixa acabaria por obrigar os bancos a recusar uma nova expansão de empréstimos em quaisquer termos. Mas antes que esse estágio seja alcançado, a alta dos preços é interrompida pelas consequências de suas próprias desigualdades inevitáveis. Essas desigualdades tornam-se mais gritantes quanto mais alto o nível geral é forçado; depois de um tempo, elas ameaçam uma séria redução de lucros para certas empresas, e os problemas dessas vítimas dissolvem aquela confiança na segurança dos créditos com a qual toda a imponente estrutura da prosperidade foi cimentada.

Quais são, então, as linhas de negócio em que os preços de venda não podem ser elevados o suficiente para evitar uma redução de lucros? Em certas linhas, os preços de venda são estereotipados pela lei, por comissões públicas, por contratos de longo prazo, pelo costume ou pela política comercial, e nestas nenhum avanço, ou apenas avanços módicos, podem ser feitos. Em outras linhas, os preços estão sempre sujeitos às chances incalculáveis das colheitas; nestas, o valor de mercado de todos os estoques acumulados de materiais e bens acabados oscila com os relatórios de safra. Em algumas linhas, a construção recente de novos equipamentos aumentou a capacidade de produção mais rápido do que a demanda por suas mercadorias se expandiu sob a influência repressiva dos altos preços que devem ser cobrados para evitar uma redução nos lucros. A falta de disposição dos investidores em firmar novos contratos ameaça prejuízo não apenas para as firmas contratantes de todos os tipos, mas também para todas as empresas das quais elas compram materiais e suprimentos. As altas taxas de juros não apenas freiam a demanda atual por mercadorias de vários tipos, mas também obstruem o esforço de manter os preços, mantendo grandes estoques de mercadorias fora do mercado até que possam ser vendidos com melhor vantagem. Finalmente, o próprio sucesso de outras empresas em elevar os preços de venda rápido o suficiente para defender seus lucros agrava as dificuldades dos homens que estão em apuros. Pois para estes últimos, cada novo aumento nos preços dos produtos que compram significa uma nova tensão sobre seus recursos já esgotados.

À medida que a prosperidade se aproxima de seu auge, desenvolve-se um forte contraste entre as perspectivas de negócios de diferentes empresas. Muitas, provavelmente a maioria, estão ganhando mais dinheiro do que em qualquer estágio anterior do ciclo econômico. Mas uma minoria importante, pelo menos, enfrenta a perspectiva de lucros declinantes. Quanto mais intensa se torna a prosperidade, maior cresce esse grupo ameaçado. É apenas uma questão de tempo para que essas condições, geradas pela prosperidade, forcem algum ajustamento radical.

Ora, tal declínio nos lucros ameaça consequências mais graves do que a falha em realizar dividendos esperados. Pois desperta dúvidas sobre a segurança dos créditos pendentes. O crédito comercial baseia-se primordialmente no valor capitalizado dos lucros presentes e prospectivos, e o crédito pendente no zênite da prosperidade é ajustado às grandes expectativas que prevalecem quando o comércio é enorme, os preços são altos e os homens de negócios otimistas. O aumento nas taxas de juros já estreitou as margens de segurança por trás dos créditos, reduzindo o valor capitalizado de dados lucros. Quando os próprios lucros começam a vacilar, a perspectiva piora. Credores cautelosos temem que o encolhimento na avaliação de mercado das empresas que lhes devem dinheiro não deixe segurança adequada para o reembolso. Assim, eles começam a recusar renovações de empréstimos antigos às empresas que não podem evitar um declínio nos lucros e a pressionar por uma liquidação das contas pendentes.

Assim, a prosperidade acaba por trazer condições que iniciam uma liquidação dos enormes créditos que ela acumulou. E no curso dessa liquidação, a prosperidade funde-se em crise.

3. CRISES

Uma vez iniciado, o processo de liquidação estende-se rapidamente, em parte porque a maioria das empresas que são chamadas a liquidar suas obrigações vencidas, por sua vez, exerce pressão semelhante sobre seus próprios devedores, e em parte porque, apesar de todos os esforços para manter em segredo o que está acontecendo, a notícia acaba vazando e outros credores alarmam-se.

Enquanto este ajustamento financeiro está em progresso, o problema de obter lucros em transações correntes é subordinado ao problema mais vital de manter a solvência. Os gerentes de negócios concentram suas energias em prover para seus passivos pendentes e em zelar por seus recursos financeiros, em vez de impulsionar suas vendas. Em consequência, o volume de novos pedidos cai rapidamente. Isto é, os fatores que já estavam obscurecendo as perspectivas de lucros em certas linhas de negócio são reforçados e estendidos. Mesmo que a esmagadora maioria das empresas consiga atender à demanda de pagamento, o teor dos negócios muda. A expansão dá lugar à contração, embora sem um solavanco violento. As taxas de desconto sobem mais do que o normal, os títulos e as mercadorias caem de preço e, à medida que os pedidos antigos são concluídos, as forças de trabalho são reduzidas; mas não há uma epidemia de falências, nem corrida aos bancos, e nenhuma interrupção espasmódica dos processos comerciais ordinários.

No extremo oposto das crises desta ordem branda estão as crises que degeneram em pânicos. Quando o processo de liquidação atinge um elo fraco na cadeia de créditos interligados e a falência de alguma empresa conspícua espalha um alarme irracional entre o público de negócios, então os bancos são subitamente forçados a enfrentar uma dupla pressão — um aumento acentuado na demanda por empréstimos e na demanda por reembolso de depósitos. Se os bancos provarem ser capazes de honrar ambas as demandas sem hesitar, o alarme rapidamente desaparece. Mas se, como aconteceu duas vezes na América desde 1890, muitos empresários solventes tiverem acomodação recusada a qualquer preço, e se os depositantes tiverem o pagamento integral recusado, o alarme transforma-se em pânico. Uma restrição de pagamentos pelos bancos dá origem a um ágio sobre a moeda corrente, ao entesouramento de dinheiro vivo e ao uso de vários substitutos ilegais para o dinheiro. Uma recusa dos bancos em expandir seus empréstimos, e mais ainda uma política de contração, envia as taxas de juros para três ou quatro vezes os seus valores habituais e causa suspensões forçadas e falências. O governo é solicitado para auxílio extraordinário, esforços frenéticos são feitos para importar ouro, certificados de empréstimo de câmaras de compensação são emitidos e a circulação de notas bancárias aumenta tão rapidamente quanto é possível sob o sistema existente. As cobranças atrasam, as taxas de câmbio doméstico são desestruturadas, os trabalhadores são demitidos porque os empregadores não conseguem dinheiro para as folhas de pagamento ou temem que não serão pagos pelas mercadorias quando entregues, os estoques caem para níveis extremamente baixos, mesmo os melhores títulos declinam um pouco de preço, os mercados de mercadorias são desorganizados por vendas de sacrifício e o comércio é violentamente contraído.

O fato de que ocasionalmente crises ainda degenerem em pânicos na América, mas não na Grã-Bretanha, França ou Alemanha, decorre primordialmente de diferenças na organização e prática bancária. Em cada um dos três países europeus, a prevalência do sistema de agências bancárias e a existência de um banco central organiza o sistema bancário como um todo, de modo que as reservas podem ser aplicadas onde e quando necessário, embora constituam apenas uma pequena porcentagem do total de passivos à vista de todas as agências. Em contraste acentuado com a política dos bancos americanos, o banco central não apenas mantém uma reserva muito acima dos requisitos imediatos em tempos normais, mas também a usa corajosamente em tempos de estresse. Como resultado, os empresários europeus não precisam temer nem uma recusa em emprestar, nem uma restrição de pagamentos pelos bancos dos quais dependem. E quando o depositante pode obter seu dinheiro em caso de necessidade e o empresário solvente pode tomar emprestado, há pouca chance de pânico.

4. DEPRESSÃO

O encerramento de um pânico é geralmente seguido pela reabertura de numerosas empresas que foram fechadas durante as semanas de pressão mais severa. Mas baseado, como é, principalmente na finalização de pedidos recebidos mas não completamente executados durante a prosperidade precedente, ou no esforço de processar e comercializar grandes estoques de materiais já em mãos ou contratados, este pronto reaquecimento da atividade é parcial e de curta duração. Ele chega ao fim à medida que este trabalho é gradualmente concluído, porque novos pedidos não estão surgindo em número suficiente para manter as usinas e fábricas ocupadas.

Segue-se um período durante o qual a depressão se espalha por todo o campo dos negócios e torna-se mais severa. A demanda dos consumidores declina em consequência de demissões em massa de assalariados, da exaustão gradual das economias e de reduções em outras classes de renda familiar. Com a demanda dos consumidores, cai a demanda das empresas por matérias-primas, suprimentos correntes e equipamentos usados na fabricação de bens de consumo. Ainda mais severo é o encolhimento na demanda dos investidores por obras de construção de todos os tipos, uma vez que poucos indivíduos ou empresas desejam investir dinheiro em novos empreendimentos de negócios enquanto o comércio permanecer deprimido e o nível de preços estiver em declínio. A contração dos negócios que resulta desses vários encolhimentos na demanda é cumulativa, uma vez que cada redução no emprego causa uma redução na demanda dos consumidores, e cada declínio na demanda dos consumidores deprime a demanda atual das empresas e desencoraja o investimento, causando assim novas demissões de funcionários e reduzindo a demanda dos consumidores mais uma vez.

Com a contração do comércio vai uma queda nos preços. Pois, quando os pedidos correntes são insuficientes para empregar o equipamento existente, a concorrência pelo negócio que houver torna-se mais acirrada. Este declínio espalha-se através dos canais comerciais regulares que conectam uma empresa a outra e é cumulativo, uma vez que cada redução de preço facilita, se não forçar, reduções em outros preços, e estas últimas reduções reagem, por sua vez, para causar novas reduções no ponto de partida.

Assim como a alta nos preços que acompanhou a retomada, a queda que acompanha a depressão é caracterizada por certas diferenças de grau que ocorrem regularmente. Os preços de atacado caem mais rápido que os de varejo, os preços de bens de produção mais rápido que os de bens de consumo, e os preços de matérias-primas mais rápido que os de produtos manufaturados. Os preços de produtos minerais brutos seguem um curso mais regular do que os de produtos florestais, agrícolas ou animais brutos. Comparados com os números de índices gerais de preços de mercadorias no atacado, os números de índices de salários e juros em empréstimos de longo prazo declinam menos, enquanto os números de índices de taxas de desconto e de ações declinam mais. O único grupo importante de preços a subir diante da depressão é o dos títulos de alta qualidade.

Naturalmente, a contração do comércio e a queda nos preços reduzem a margem de lucros presentes e prospectivos, espalham o desânimo entre os empresários e freiam o empreendedorismo. Mas eles também põem em movimento certos processos de ajustamento pelos quais a depressão é gradualmente superada.

Os custos primários de fazer negócios são reduzidos pela queda rápida nos preços das matérias-primas e dos empréstimos bancários, pelo aumento acentuado na eficiência da mão de obra que ocorre quando o emprego é escasso e os homens estão ansiosos por manter seus empregos, e por uma economia mais rigorosa por parte dos gerentes. Os custos suplementares também são reduzidos pela reorganização de empresas que realmente se tornaram ou ameaçam tornar-se insolventes, pela venda de outras empresas a valores baixos, pela redução em aluguéis e refinanciamento de empréstimos, pela baixa de dívidas ruins e depreciação de propriedades, e por admitir que uma recapitalização de empresas comerciais — correspondente aos preços mais baixos das ações — foi efetuada com base em lucros menores.

Enquanto os custos ainda estão sendo reduzidos, a demanda por mercadorias para de encolher e começa a expandir-se lentamente — uma mudança que geralmente ocorre no segundo ou terceiro ano de depressão. Estoques acumulados deixados pela prosperidade são gradualmente esgotados e o consumo corrente exige produção corrente. Roupas, móveis, máquinas e outros artigos moderadamente duráveis que foram usados o máximo possível são finalmente descartados e substituídos. A população continua a aumentar a uma taxa bastante uniforme: as novas bocas devem ser alimentadas e as novas costas vestidas. Novos gostos aparecem entre os consumidores e novos métodos entre os produtores, dando origem à demanda por produtos inovadores. O mais importante de tudo é que a demanda de investimento por equipamentos industriais revive; pois, embora a poupança possa diminuir, ela não cessa; com a cessação das vendas judiciais e reorganizações corporativas, as oportunidades de comprar empresas antigas a preços de pechincha tornam-se menos frequentes; os capitalistas tornam-se menos tímidos à medida que a crise recua no passado; as baixas taxas de juros em títulos de longo prazo incentivam o endividamento; os aperfeiçoamentos técnicos acumulados de vários anos podem ser utilizados; e os contratos podem ser firmados em condições altamente favoráveis quanto ao custo e execução imediata.

Uma vez que essas várias forças puseram o comércio a expandir-se novamente, o aumento prova-se cumulativo, embora por um tempo o ritmo de crescimento seja mantido lento pela queda contínua dos preços. Mas enquanto esta última mantém a pressão sobre os empresários e impede que o aumento nos pedidos produza uma alta rápida nos lucros, as perspectivas de negócios tornam-se gradualmente mais brilhantes. Dívidas antigas foram pagas, estoques acumulados de mercadorias absorvidos, empresas fracas reorganizadas, os bancos estão fortes — todas as nuvens no horizonte financeiro desapareceram. Tudo está pronto para uma retomada da atividade, que começará sempre que alguma circunstância afortunada der um impulso repentino à demanda ou, na ausência de tal evento, quando o lento crescimento dos negócios tiver preenchido os livros de pedidos e pavimentado o caminho para uma nova alta nos preços. Tal é o estágio do ciclo econômico com o qual a análise começou e, tendo explicado seu próprio início, a análise termina.

NOTA
RELAÇÃO DA TEORIA PRECEDENTE DOS CICLOS ECONÔMICOS COM OUTRAS TEORIAS CORRENTES EM 1913

Quando considerada em relação à análise precedente, nenhuma das teorias dos ciclos econômicos resumidas na Parte I, Capítulo I, parece ser demonstravelmente errada, mas nem uma delas parece ser inteiramente adequada. Talvez um esforço para sugerir as limitações de cada teoria possa lançar luz adicional sobre um problema que é obscuro, na melhor das hipóteses — certamente tal esforço revelará o quanto a presente discussão tomou emprestado de suas antecessoras.

A teoria de Beveridge, de que o esforço competitivo de cada produtor para abocanhar o máximo possível do mercado para suas próprias mercadorias leva necessariamente a excessos de oferta, coloca toda a ênfase em um único aspecto em um conjunto de processos que se mostraram excessivamente complexos. Mas, para fazer justiça a Beveridge, ele está longe de professar avançar uma teoria adequada das crises. E o único fator que ele enfatiza certamente conta no sentido de que todo o curso dos ciclos econômicos seria profundamente alterado se a competição entre produtores fosse suprimida. Mas as alegações feitas livremente há algumas décadas, tanto na América quanto na Alemanha, de que o desenvolvimento de trustes e cartéis poria fim às crises, não foram sustentadas pela história dos negócios. As indústrias nas quais a combinação fez maiores progressos no sentido de regular a competição ainda sentem a tensão causada pelo aumento do custo da mão de obra e pela tensão nos mercados de investimento e monetário. Nem os gerentes dessas grandes combinações foram capazes de evitar a redução nos lucros causada pelo declínio da demanda por seus produtos quando a prosperidade passou de seu zênite.

As várias formas da teoria do subconsumo já foram criticadas. Seu valor reside em destacar um dos possíveis obstáculos para prevenir invasões de custos sobre os lucros. Mas nunca foi demonstrado que a demanda dos consumidores fica atrás da oferta antes de uma crise ter começado. E quaisquer que sejam os fatos, há evidências de que as tensões que realmente precipitaram as crises dos anos recentes apareceram mais cedo em outros ramos do comércio do que naqueles que atendem às necessidades dos consumidores (ver Cap. 2, iv, 2, D).

A teoria de Spiethoff sobre a produção desequilibrada de equipamentos industriais e de bens complementares fornece dois elementos na discussão precedente. Mas aqui esses elementos foram colocados em um cenário um pouco diferente. A subprodução de bens complementares aparece sob a forma de um aumento nos custos de materiais, mão de obra e afins, tão rápido que ameaça reduzir os lucros prospectivos. A superprodução de equipamentos industriais é apresentada como uma das razões pelas quais é difícil avançar os preços de venda suficientemente para compensar esse aumento nos custos. Se este cenário ajuda a tornar mais clara a importância da análise de Spiethoff, ele também serve para indicar que outros elementos além daqueles sobre os quais ele se alonga devem ser levados em conta.

A teoria de Hull de que o alto custo das obras de construção em períodos de prosperidade intensa causa um declínio severo na demanda de investimento e, portanto, uma crise, que se espalha de indústrias como a do aço para todos os ramos de negócios, foi incorporada na análise precedente. Mas ali ela é combinada com inúmeras outras ideias que o Sr. Hull negligenciou.

A teoria de Lescure sobre variações nos lucros prospectivos combinada com a teoria de Veblen sobre a discrepância entre os lucros prospectivos e a capitalização corrente fornece a estrutura geral do Capítulo 2. Mas a análise desses dois autores foi desenvolvida com o auxílio de sugestões extraídas de outras fontes, e foi feita uma tentativa de testar certas de suas suposições pelo uso de materiais estatísticos.

Entre as sugestões assim utilizadas estão as de Sombart, Carver e Irving Fisher. Embora todas as três sugestões pertençam a uma análise de como a prosperidade traz um declínio nos lucros prospectivos e, assim, mina a base do crédito, nenhuma delas pode isoladamente sustentar-se adequadamente.

Tensões dentro do sistema de preços podem muito bem surgir, como sustenta Sombart, porque a oferta de bens orgânicos não acompanha o mesmo ritmo que a oferta de minerais e afins; mas, a menos que a análise precedente seja pura imaginação, outros fatores estão em jogo, mais regulares em sua operação do que as incertezas das colheitas, e não menos potentes em invadir os lucros.

A teoria de Carver sobre as flutuações de preços dessemelhantes entre bens de produção e de consumo é outra forma de explicar que a prosperidade traz um rápido aumento de equipamento industrial e uma crise ocorre quando a crescente oferta de bens não pode ser vendida com lucro. Naturalmente, este aspecto foi considerado; mas está longe de ser o único — uma conclusão que o Sr. Carver estaria pronto a reconhecer se estivesse lidando com o problema em uma escala mais ampliada.

Da mesma forma, o ajuste retardado das taxas de juros às flutuações no nível de preços, utilizado pelo Sr. Fisher para explicar a ocorrência de crises, é apenas um entre vários fatores que alargam as margens de lucro nos estágios iniciais da prosperidade e as estreitam em um estágio posterior.

Finalmente, a teoria do Sr. Johannsen sobre a "poupança prejudicada" foi usada incidentalmente: isto é, no Capítulo 4 reconhece-se que investimentos em títulos de refinanciamento, em títulos de empresas sendo reorganizadas e em propriedades vendidas judicialmente, não têm tal efeito estimulante sobre os negócios quanto os investimentos em novo equipamento industrial. É em parte — mas apenas em parte — porque investimentos do primeiro tipo predominam nos meses que sucedem uma crise que a produção corrente cai.

Enquanto as provas deste capítulo estavam em minhas mãos, apareceu o esboço preliminar da projetada teoria das crises da Sra. Minnie Throop England ("Economic Crises", Journal of Political Economy, abril de 1913). A Sra. England parece ter formulado o problema de forma muito semelhante à que eu fiz, ter empregado métodos parecidos e ter chegado a conclusões amplamente semelhantes.

II. Diversidades entre os Ciclos Econômicos e Suas Causas

Qualquer análise que trace o curso geral dos processos que trazem prosperidade, crise e depressão inevitavelmente deixa uma impressão enganosa de uniformidade entre os ciclos econômicos. Na verdade, esses ciclos diferem amplamente em duração, em intensidade, na proeminência relativa de seus vários fenômenos e na sequência de suas fases.

1. AS DIVERSIDADES

Se os anos entre uma crise e a seguinte forem tomados como a duração de um ciclo econômico, os ciclos britânico, francês e alemão iniciados em 1890 duraram dez anos, e aqueles iniciados em 1900 duraram sete anos. Os ciclos americanos contemporâneos têm variações mais amplas: três anos de 1890 a 1893; dez anos de 1893 a 1903; e quatro anos de 1903 a 1907. Em vista dessas diversidades, a noção de que as crises têm um período regular de recorrência é claramente equivocada.

A experiência americana em 1893-1903 refuta a ideia de que os ciclos econômicos consistem sempre em uma prosperidade inabalável, seguida por uma crise bem definida e uma depressão ininterrupta. A depressão foi interrompida pela retomada em 1895 e, em seguida, agravada pela extrema escassez monetária de 1896. A retomada que começou em 1897 foi freada por um curto período pelo início da guerra com a Espanha na primavera de 1898. O período próspero seguinte foi prejudicado pelos problemas da bolsa de valores de 1899 e 1901 e pelo abrandamento temporário dos negócios gerais em 1900.

Como exemplos de diferenças de intensidade, o contraste entre os ciclos francês e americano é o mais marcante. Mas mesmo dentro dos limites de um país, ciclos sucessivos trazem crises relativamente brandas como as de 1890 e 1903 na América e pânicos severos como os de 1893 e 1907; períodos de depressão relativamente breves e moderados como os de 1891 e 1904, e períodos relativamente longos e drásticos como o de 1894-96; períodos de prosperidade relativamente transitórios como os de 1892 e 1909, e períodos relativamente longos de prosperidade intensa como os de 1898-1902 e 1905-07.

Finalmente, não há dois períodos de prosperidade, crise ou depressão que tenham exatamente a mesma combinação de elementos. A especulação é por vezes desenfreada durante a prosperidade, como na América em 1901, e por vezes mantida firmemente sob controle, como na Grã-Bretanha em 1906-07. Uma contração de crédito é frequentemente a característica mais conspícua das crises, como nos pânicos americanos de 1893 e 1907; enquanto em outras vezes ela desempenha um papel menor em comparação com o volume declinante de novos negócios, como na crise britânica de 1907. Diferentes ramos do comércio e diferentes seções do país revelam-se como as principais sedes de atividade, as principais fontes de tensão e os que mais sofrem com a depressão em ciclos sucessivos. De fato, as dissimilaridades deste tipo são tão numerosas e tão óbvias que não vale a pena citar mais do que os poucos exemplos dados na discussão anterior.

2. SUAS CAUSAS

Muitas dessas divergências entre os ciclos econômicos devem-se a eventos que surgem de outras fontes que não os negócios. Pois o mecanismo da economia monetária é tão delicado que as perspectivas de lucro de alguém são afetadas pelas notícias de cada dia. O mais importante de todos esses fatores estranhos a longo prazo são as chances do clima que tornam as colheitas boas ou ruins, e assim afetam os preços dos produtos agrícolas, o poder de compra das comunidades agrícolas, os ganhos das ferrovias, etc., etc. A eclosão da guerra ou a celebração da paz, distúrbios de ordem interna, terremotos, conflagrações, epidemias, mudanças nos padrões monetários, revisões tarifárias, políticas governamentais relativas a corporações, alterações na produção de ouro, aperfeiçoamentos na técnica industrial, a abertura de novas terras para colonização, o esgotamento de recursos naturais, a mudança de rotas comerciais — estas e mil outras coisas dificilmente deixarão de ajudar ou prejudicar algum empreendimento comercial. Se o círculo que elas atingem for grande e seus efeitos pronunciados, elas sem dúvida dão um toque peculiar ao ciclo econômico em que caem. Particularmente em tempos de transição de uma fase do ciclo para outra, quando o rumo da corrente é incerto, a influência de tais eventos é frequentemente acentuada. Mas no meio da prosperidade, crise ou depressão, muitas vezes parece que a comunidade de negócios não dá atenção a notícias que não estejam de acordo com o teor da época.

Menos óbvias, mas mais persistentes em seu efeito sobre o curso e o caráter dos ciclos econômicos, são as mudanças que ocorrem continuamente na organização e prática de negócios e na importância relativa das diferentes indústrias. Por exemplo, tanto a extravagância dos auges quanto a violência dos pânicos foram temperadas por uma organização mais estreita, conhecimento mais amplo e políticas mais firmes entre os bancos. Novamente, o rápido desenvolvimento da manufatura e o declínio da construção ferroviária em importância entre as indústrias americanas ajudaram a tornar os ciclos econômicos de 1900-10 diferentes daqueles de 1880-90. Mais uma vez, o crescente domínio de grandes corporações que dependem do público investidor para seus fundos dá maior influência aos mercados de títulos e ações do que esses mercados exerciam nos dias das empresas familiares. O amplo contraste entre os ciclos francês e americano mostra quão poderosamente o desenvolvimento relativo da poupança ou do empreendimento afeta a intensidade tanto da prosperidade quanto da depressão. Outras mudanças que reagem sobre os ciclos econômicos são a extensão do controle monopolista, a integração da indústria, a organização do trabalho com sua padronização de taxas salariais e, em geral, o reajuste dos negócios para atender a mudanças no ambiente material, político ou social.

A reação de tais mudanças contemporâneas na organização econômica sobre o caráter de ciclos econômicos sucessivos pode ser difícil de traçar. Mas cada nova mudança realizada torna-se a base sobre a qual novas mudanças procedem. Ou seja, as amplas mudanças da organização econômica são cumulativas, como as mudanças menores que fazem com que cada fase de cada ciclo econômico evolua para a sua sucessora. E, sendo cumulativas, sua influência dominante sobre os fenômenos dos ciclos econômicos destaca-se claramente no decorrer dos anos. Por isso, os economistas de cada geração provavelmente verão razões para reformular a teoria dos ciclos econômicos que aprenderam em sua juventude.

III. Os Ciclos Econômicos na História Econômica

1. GÊNESE DOS CICLOS ECONÔMICOS

Se o termo crise econômica for feito para cobrir qualquer distúrbio sério dos processos usuais de produção e distribuição de bens, então tais crises são pelo menos tão antigas quanto os registros econômicos. "Atos de Deus" destrutivos e a ira desperdiçadora do homem frequentemente trouxeram angústia aguda sobre comunidades de pessoas industriosas muito antes que a organização econômica assumisse sua forma atual e os ciclos econômicos começassem a correr em um curso regular.

Mais modernos em seu caráter foram os distúrbios que se seguiram à organização do crédito. No século XVI, por exemplo, as dívidas públicas da França, Espanha, Áustria e outros países haviam crescido tanto que os repúdios reais tornaram-se uma fonte prolífica de crises financeiras em centros como Antuérpia e Lyon. O surgimento do sistema bancário expôs os círculos mercantis a outros perigos ainda. Os ourives-banqueiros em Londres, para citar instâncias específicas, foram submetidos a corridas e seus clientes a pânicos quando os holandeses queimaram a frota inglesa em Chatham, no Tâmisa, em 1667, e novamente em 1672, quando Carlos II interrompeu os pagamentos do Tesouro. Mas, embora tais episódios possam ser chamados justamente de crises financeiras, eles diferem de suas contrapartidas recentes tanto por afetarem apenas alguns ramos do comércio quanto por dependerem diretamente da guerra ou dos embaraços fiscais do governo.

Mesmo no século XVIII, a maioria das crises inglesas surgiu de fontes que não os negócios — embora nesta época a Inglaterra já tivesse assumido definitivamente a liderança na organização econômica. Em 1708, os ourives-banqueiros aproveitaram-se de relatos de conspirações jacobitas e de preparativos na França para uma incursão na Escócia para atacar o crédito do Banco da Inglaterra. Em 1720, a "Bolha dos Mares do Sul" estourou, encerrando uma mania quase inacreditável de especulação em ações. A crise seguinte ocorreu em 1745, quando o Pretendente com seus Highlanders penetrou a 120 milhas de Londres. O fim da Guerra dos Sete Anos foi seguido em 1763 por uma viva especulação e colapso na Bolsa de Valores. Novamente em 1771-73, a celebração da paz levou à especulação e a um colapso. Cinco ou seis anos depois, perdas trazidas pela guerra com as Colônias Americanas causaram sérias dificuldades comerciais. Quando esta guerra terminou, em 1783, a paz mais uma vez deu origem a uma expansão repentina de negócios e a expansão terminou em uma crise. Finalmente, em 1793, veio o que o historiador desses eventos chama de a primeira das grandes crises industriais da Inglaterra, seguida por depressão nos negócios em geral.

Este breve relato indica que os ciclos econômicos são muito posteriores em seu aparecimento do que as crises econômicas, ou mesmo estritamente financeiras. Na própria Inglaterra, eles parecem não ter começado antes do final do século XVIII. Mas quando apareceram, foi na forma de uma extensão sobre todos os ramos da indústria de dificuldades não muito diferentes daquelas que vinham sendo sofridas por mais de cem anos por grandes capitalistas, banqueiros e especuladores em ações. Com esta extensão no escopo, veio uma mudança na importância relativa das causas. No passado, o enfraquecimento do crédito costumava ser causado pela guerra, pela celebração da paz ou por alguma violação das obrigações financeiras por parte do governo. No futuro, passaria a ser causado com mais frequência por tensões geradas dentro do próprio mundo dos negócios.

A razão para ambas as mudanças residia na extensão gradual da empresa comercial altamente organizada de seus centros iniciais de comércio exterior, mineração, finanças e bancos sobre o amplo campo da manufatura e do comércio doméstico — uma extensão que acompanhou a Revolução Industrial. À medida que o artesanato deu lugar a fábricas geridas por princípios de negócios, atendendo a um mercado amplo e incerto, entrando livremente em contratos de longo prazo, exigindo um pesado investimento de capital fixo e usando dinheiro emprestado em escala liberal, o círculo de empresas afetadas por dificuldades financeiras cresceu constantemente, e o perigo de que dificuldades financeiras surgissem da conduta dos assuntos de negócios tornou-se cada vez maior.

Além disso, à medida que a manufatura e o comércio doméstico tornaram-se dependentes dos mercados monetário e de investimento, as crises começaram regularmente a ser seguidas por períodos de liquidação nos círculos manufatureiros e mercantis. Mas como as dificuldades financeiras agora surgiam em grande parte do embaraço de fabricantes e comerciantes, uma liquidação completa por parte destes últimos encerrava as condições que haviam causado a crise e pavimentava o caminho para uma retomada da atividade econômica. E como tanto a acumulação quanto o relaxamento dessas tensões ocorriam regularmente como um subproduto dos processos de negócios, as crises irregulares e imprevisíveis do passado transformaram-se gradualmente nos ciclos de prosperidade, crise e depressão com os quais os homens passaram a contar e que estão começando a prever.

À medida que a Revolução Industrial e suas mudanças concomitantes na organização do comércio e do transporte se espalharam para outros países, estes últimos começaram a desenvolver os fenômenos dos ciclos econômicos já familiares na Inglaterra. E à medida que a empresa comercial completava seu domínio sobre a manufatura, o comércio atacadista, as finanças, os transportes, a mineração e a exploração madeireira, e começava a invadir o comércio varejista, as profissões e a agricultura, as oscilações cíclicas de expansão e contração trouxeram cada vez mais homens sob seu domínio imediato.

2. O DOMÍNIO DO HOMEM SOBRE O FUNCIONAMENTO DA ECONOMIA MONETÁRIA

Os ciclos econômicos, portanto, aparecem naquele estágio da história econômica em que o processo de fabricação e distribuição de bens é organizado principalmente na forma de empresas comerciais conduzidas para o lucro. Esta forma de organização econômica foi desenvolvida gradualmente a partir de formas anteriores por sucessivas gerações de homens que pensaram em obter alguma vantagem de cada passo sucessivo. Mas o maquinário complicado e o funcionamento do sistema não são totalmente dominados nem mesmo pela geração atual de empresários e, recorrentemente, o maquinário financeiro inflige sofrimento grave sobre nós que o usamos. Porque não aprendemos como evitar que os custos invadam os lucros e que a escassez se acumule nos mercados monetários, como manter estável a construção de novo equipamento industrial, como controlar a capitalização de mercado das empresas e como evitar expansões e contrações espasmódicas de créditos — porque nosso conhecimento teórico e nossa habilidade prática são deficientes em relação a esses assuntos técnicos, não podemos manter a prosperidade por mais do que alguns anos de cada vez.

No entanto, no último século, progredimos em direção ao domínio sobre os processos da economia monetária. A Mania das Tulipas na Holanda, o Esquema dos Mares do Sul na Inglaterra e a Bolha do Mississippi na França não têm rivais dignos nas décadas recentes. Mesmo a excitação especulativa que precedeu a crise de 1873 nos estados alemães e na América raramente foi igualada desde 1890. Por uma combinação de várias agências, como a regulamentação pública dos prospectos de novas empresas, legislação apoiada por administração eficiente contra promoções fraudulentas, requisitos mais rígidos por parte das bolsas de valores relativos aos títulos admitidos nas listas oficiais, agências mais eficientes para dar informações aos investidores e uma política mais conservadora por parte dos bancos em relação a auges especulativos, aprendemos a evitar certos dos erros mais imprudentes cometidos por gerações anteriores. Novamente, por meio de dura experiência, os bancos europeus, pelo menos, aprenderam métodos de controlar uma crise e evitar que ela se transforme em pânico. A "integração da indústria" também fez algo, embora menos do que é frequentemente alegado, no sentido de estabilizar o curso dos negócios, tanto concentrando o poder nas mãos de funcionários experientes quanto moderando as flutuações extremas de preços.

3. PROPOSTAS PARA CONTROLAR OS CICLOS ECONÔMICOS

O que já foi realizado nessas direções para controlar nossa maquinaria de negócios pode muito bem ser o prenúncio de maiores conquistas. Três linhas promissoras de esforço são sugeridas pelas propostas para reorganizar o sistema bancário americano, para usar as compras governamentais e ferroviárias como um volante de equilíbrio dos negócios e para "estabilizar o dólar".

A primeira proposta visa primordialmente evitar que as crises se tornem pânicos. Sua viabilidade foi demonstrada pela experiência europeia e canadense, e planos práticos foram elaborados em detalhes pela Comissão Monetária, comitês do congresso, associações de banqueiros e investigadores privados. De fato, a reforma bancária é tanto a mais necessária quanto a mais fácil entre todas as mudanças que prometem aumentar nosso controle sobre o funcionamento da economia monetária. Mas a literatura sobre este assunto é tão vasta e acessível que não há necessidade de entrar em detalhes neste livro.

A segunda proposta, de usar as compras governamentais e ferroviárias como um volante para estabilizar o mecanismo de negócios, visa primordialmente mitigar a severidade das depressões. Ela foi formulada de forma mais definitiva na França e na Inglaterra. Em 1907, o Ministro de Obras Públicas da França determinou que fosse feita uma investigação sobre o efeito das crises sobre as ferrovias. O relatório argumentou que é viável para os grandes sistemas ferroviários distribuir seus pedidos de material rodante, e assim por diante, sistematicamente ao longo de todo o período de um ciclo econômico, de forma a reduzir os pedidos agora feitos em anos de grande atividade e aumentar o volume nos anos de marasmo. Esta mudança seria para a vantagem financeira das ferrovias, na medida em que o equipamento é mais barato para construir em anos de depressão; para a vantagem dos expedidores, na medida em que uma provisão mais oportuna de material rodante diminuiria a frequência de bloqueios de carga e escassez de vagões; e para a vantagem do público, na medida em que o emprego seria estabilizado em uma indústria importante. Como forma de aplicar estas ideias, o Ministro de Obras Públicas, em maio de 1907, convidou as ferrovias a submeterem um programa definido para a compra de material rodante abrangendo os anos 1907-10.

Mais ambicioso em escopo e diferente em ênfase é o esquema elaborado pelo Sr. e Sra. Sidney Webb. Como parte de um plano abrangente para a prevenção da indigência, eles propõem que uma parte dos contratos governamentais seja retida em anos de atividade intensa e liberada em anos de baixa. Os números do Dr. A. L. Bowley, nos quais eles confiam, indicam que se apenas 3 ou 4 por cento dos pedidos do governo fossem tratados desta forma, uma grande parte do desemprego agora causado pelas depressões cíclicas poderia ser contrabalançada. Entre os itens que poderiam ser relegados a um programa de dez anos e contratados apenas quando o comércio desse sinais de queda, eles mencionam: "...metade da dotação anual para reconstrução e multiplicação de edifícios governamentais...; metade da provisão anual normal para suprimentos como cobertores, lona e pano cáqui, dos quais há sempre um grande estoque; o todo (ou metade) da soma alocada anualmente para a colocação gradual de fios telegráficos subterrâneos e a extensão gradual do telefone em cada pequena aldeia; todo o trabalho de impressão, como os relatórios da Comissão de Manuscritos Históricos e a história oficial da Guerra da África do Sul; pelo menos metade da despesa anual no desenvolvimento das florestas governamentais...; uma proporção considerável dos subsídios do Conselho de Educação para a construção de novas faculdades de formação e escolas secundárias; alguma parte da construção naval normal do ano...; pelo menos metade da dotação anual para novos estandes de tiro e salas de treinamento para a Força Territorial; a maior parte das despesas de capital do Conselho de Distritos Congestionados na Irlanda, e assim por diante. E a isto deve ser adicionado o total das somas, totalizando mais de um milhão por ano, já colocadas à disposição dos Comissários de Desenvolvimento e do Conselho Rodoviário. Está muito claro [concluem eles] que existe, no conjunto, um montante muito grande — do qual o total de quatro milhões por ano poderia ser facilmente selecionado — e uma variedade muito considerável de despesas que poderiam, sem qualquer inconveniente apreciável, ser reorganizadas dentro da década."

Finalmente, se o engenhoso plano do Professor Irving Fisher de "estabilizar o dólar" fosse adotado, e se ele tivesse sucesso em manter as flutuações do nível de preços dentro de limites estreitos, o curso dos ciclos econômicos tornar-se-ia sem dúvida mais uniforme do que atualmente. Pois se cada alta ou queda nos preços das mercadorias fosse freada prontamente, um dos principais fatores que agora causam mudanças nas perspectivas de lucro seria virtualmente removido. Naturalmente, ainda poderiam ocorrer mudanças desiguais nos preços relativos de matérias-primas e produtos acabados, de bens de consumo no atacado e varejo, ou de mão de obra e empréstimos, o que afetaria as margens de lucro sem perturbar em grau equivalente o número de índice geral usado como base para compensar o dólar. Os lucros também estariam sujeitos a mudanças decorrentes de variações no volume de negócios, na eficiência da mão de obra e assim por diante. Mas como a magnitude dos efeitos agora produzidos por estes fatores resulta da sua ação cumulativa em aumentar a intensidade uns dos outros, a paralisação de um fator potente tornaria os outros também menos potentes.

A viabilidade de realizar estas propostas, exceto a reforma bancária, não foi demonstrada em grande escala. Certamente todos os três planos merecem consideração cuidadosa. Eles são trazidos aqui, entretanto, meramente para sugerir formas possíveis pelas quais o controle do homem sobre o funcionamento da economia monetária pode ser aumentado.

Além de tais reformas específicas na organização econômica, o progresso reside na direção de melhorar nossas previsões das condições de negócios. Pois quando os problemas que virão são previstos, eles podem muitas vezes ser mitigados e, por vezes, evitados.

IV. Previsão das Condições de Negócios

A incerteza que acompanha as previsões decorre principalmente das imperfeições do nosso conhecimento sobre as condições de negócios no passado imediato e no presente. Pois, uma vez que os ciclos econômicos resultam de processos de mudança cumulativa, os principais fatores que moldam o amanhã são os fatores que estavam em ação ontem e estão em ação hoje.

Ora, a economia monetária permite percepções muito desiguais sobre o seu próprio funcionamento para as diferentes classes de empresários. Estas desigualdades abrem uma rica fonte de lucro para a classe favorecida, e os seus esforços para aproveitar ao máximo as suas chances podem aumentar a violência das crises. Uma forma de aumentar o controle social sobre a atividade econômica é, portanto, democratizar o conhecimento das condições de negócios atuais já possuído por alguns. Mas os caminhos e meios para este fim podem ser discutidos melhor após as desigualdades presentes de conhecimento e as suas consequências terem sido expostas.

1. OPORTUNIDADES EXCEPCIONAIS DE CERTOS FINANCIISTAS

No que diz respeito à condição de negócios de importância crucial — as flutuações nos lucros — é especialmente difícil obter informações rápidas, confiáveis e amplas. Certos homens estão tão posicionados no mundo dos negócios, entretanto, que têm uma visão ampla sobre os lucros e são tão treinados que podem tirar conclusões bastante seguras sobre a tendência das mudanças em andamento. Os grandes capitalistas que participam ativamente em muitas empresas e que têm relações pessoais íntimas com outros capitães de indústria estão nesta posição de vantagem, assim como os gerentes dos maiores bancos. Sem dúvida, a informação à disposição mesmo destes homens no centro do sistema financeiro deixa muito a desejar em termos de escopo e precisão. Sem dúvida, as inferências que tiram são distorcidas pelas suas equações pessoais e as suas previsões divergem apreciavelmente. Mas se estes cavalheiros pudessem ser induzidos a publicar francamente o que sabem sobre o presente e o que realmente esperam para o futuro próximo, é provável que o consenso das suas opiniões raramente se provaria muito errado.

No entanto, esta clarividência superior é um ativo de negócios de valor demasiado elevado para ser dado de graça. O homem de grandes recursos que pode prever a chegada de uma crise é capaz de mudar as suas participações de forma a evitar perdas que recairão sobre os incautos e de lucrar por ter fundos à mão para comprar propriedades a preços de pechincha. Não é provável que ele leve o público à sua confiança até que os seus próprios assuntos tenham sido assim organizados. Quando ele fala, pode muito bem ser com vista a influenciar a tendência do sentimento entre os pequenos empresários e investidores para sua própria vantagem. O público, portanto, tem boas razões para ler com espírito altamente crítico as entrevistas que financistas proeminentes ocasionalmente concedem.

A vantagem desfrutada por este pequeno grupo de grandes financistas não se limita a oportunidades superiores de prever mudanças que se aproximam. Em certa medida, eles podem controlar os eventos que preveem. Esta capacidade surge principalmente da crescente centralização do poder de conceder ou reter créditos. Por um lado, o surgimento da grande corporação tornou as empresas de importância estratégica dependentes dos mercados metropolitanos para empréstimos e títulos, em vez de bancos e investidores locais. Por outro lado, os grandes bancos, companhias de seguros e casas de investimento que dominam os mercados financeiros de Nova York, Londres, Paris e Berlim desenvolveram relações íntimas e podem ser controlados por algumas pequenas camarilhas de financistas. A estes homens é dada, portanto, uma grande medida de poder sobre a concessão de empréstimos bancários, o lançamento de novos títulos e os preços de ações e obrigações em circulação. Este poder eles podem usar, se quiserem, para aumentar as tensões que a prosperidade gera. Se eles retiverem grandes somas, por exemplo, reduzem as reservas dos bancos e precipitam a revisão descendente de créditos com que uma crise começa. Se eles impedirem corporações de levantar empréstimos necessários para cumprir obrigações que vencem, forçam a nomeação de administradores judiciais, derrubam o preço das ações e criam um sentimento de desconfiança que produz novas consequências próprias.

O pouco que se sabe sobre o funcionamento interno da "alta finança" indica que este poder ainda não foi exercido com eficiência implacável. Sem dúvida, muitos empresários recuariam diante da ideia de agravar deliberadamente uma crise para o seu próprio ganho. Além disso, os financistas que têm mais poder sobre o crédito estão muitas vezes fortemente interessados em empresas industriais e temem perder dividendos na depressão que se seguiria a uma crise. Um terceiro impedimento é a obsessão de que o dólar é uma medida estável de valor. Tão acostumados os empresários ficam a tratar o dólar como constante e a imputar todas as mudanças nos preços a flutuações no valor das mercadorias cotadas, que não percebem prontamente o lucro monetário a ser obtido com as mudanças no nível geral de preços. Finalmente, mesmo nos círculos mais altos das finanças, a centralização do poder ainda não foi longe o suficiente para garantir a unanimidade de ação.

Entre estes impedimentos ao esforço de agravar as flutuações das condições de negócios, pelo menos dois parecem estar perdendo a sua força. A crescente mobilidade dos investimentos está tornando mais fácil para os financistas retirarem os seus fundos de emaranhados industriais e colocá-los de tal forma que uma depressão não traga perdas sérias. E as flutuações contínuas no nível de preços estão sempre demonstrando que os dólares são unidades mutáveis, de cujas flutuações lucros podem ser realizados. É, portanto, perfeitamente possível que os financistas explorem as suas oportunidades de agravar as crises com maior energia no futuro imediato do que o fizeram no passado recente.

Se os seus esforços nesta direção se tornarem mais enérgicos, pode-se esperar que os capitalistas à moda antiga, que estão interessados primordialmente na conduta ininterrupta das operações industriais, se juntem ao grande número de empresários menores para se protegerem de ataques ao crédito. Provavelmente, a demanda por regulamentação governamental, que se dirige atualmente principalmente para a regulamentação das corporações ferroviárias e industriais, seria então estendida à regulamentação de todas as operações financeiras envolvidas na concessão de empréstimos. Mas ainda não está claro precisamente como o governo poderia intervir para evitar que grupos poderosos de financistas retivessem dinheiro quando bem entendessem, que recusassem acomodação bancária quando escolhessem, ou que declinassem subscrever novas emissões de títulos para corporações fora de favor.

De fato, o aumento da regulamentação governamental na indústria pode tornar os capitalistas mais agressivos e mais ansiosos para se interessarem primordialmente por finanças e explorarem ao máximo as suas oportunidades de manipular o crédito. Nesse caso, uma segunda linha de defesa pode ser seguida por muitos cujas perspectivas de negócios estão em risco. Se informações confiáveis sobre lucros e créditos, tais como existem agora em forma parcial e em poucas mãos, fossem coletadas e publicadas, a vantagem estratégica dos grandes financistas na previsão de crises seria materialmente reduzida. Pois então todos os interessados poderiam tomar as precauções que as suas atividades exigissem e o seu julgamento sugerisse, quando a prosperidade fosse vista como geradora de uma crise. De fato, cada um, dentro da medida permitida pelos seus recursos pecuniários, poderia tentar transformar a crise iminente em seu próprio proveito — uma condição que resultaria em poucos serem apanhados em desvantagem. O campo deixado aberto para a exploração predatória pelos financistas seria reduzido à produção de contrações repentinas de crédito sobre as quais as estatísticas de lucros e créditos não tivessem dado qualquer aviso. Ou seja, estas campanhas correriam contra a tendência geral do desenvolvimento dos negócios e seriam, portanto, muito menos destrutivas do que campanhas reforçadas pelas tensões acumuladas pela prosperidade.

Mas tudo isto diz respeito a um futuro problemático. Como as coisas se apresentam com o público de negócios atualmente?

2. BARÔMETROS DE NEGÓCIOS DISPONÍVEIS AO PÚBLICO

O homem de negócios americano que procura manter-se informado sobre a tendência das condições de negócios em geral confia nas colunas financeiras do seu jornal diário, suplementadas talvez por um ou dois semanários financeiros e um jornal comercial especializado. Os dados que ele pode compilar a partir destas fontes cobrem uma gama considerável.

Os preços das mercadorias no atacado são representados tanto por cotações reais para os grandes produtos básicos do comércio como por números de índices como o da Bradstreet. Os preços dos empréstimos à vista e a prazo de trinta dias a seis meses são relatados para Nova York, juntamente com as taxas de mercado e bancárias em Londres, Paris e Berlim. Os preços dos títulos negociados na bolsa de Nova York são publicados em detalhe e, para a tendência geral do mercado, existem registros convenientes, como os preços reais médios do Wall Street Journal para vinte ações ferroviárias e doze industriais.

As flutuações nos negócios devem ser estimadas a partir de várias fontes: compensações bancárias, receitas brutas das ferrovias, número de vagões parados, importações e exportações, produção de carvão, cobre, ferro gusa e aço, embarques de cereais, algodão, gado, etc. Os relatórios governamentais de safras ajudam a prever o provável estado do comércio em várias seções agrícolas. Menos sistemáticas, mas muitas vezes úteis, são as numerosas análises de condições de negócios em diferentes ramos e cidades, publicadas em intervalos regulares por vários jornais semanais.

Informações sobre a moeda são fornecidas pelas estimativas oficiais do estoque monetário, por relatórios de importações e exportações de ouro, pelos movimentos registrados de dinheiro para dentro e para fora dos bancos de Nova York, e pelos números relativos à produção e ao consumo industrial de ouro, e a distribuição de dinheiro entre os bancos e o público — embora estas últimas três declarações surjam apenas uma vez por ano. Relativamente aos bancos, existem declarações telegráficas das instituições centrais da Europa, relatórios semanais das câmaras de compensação de Nova York, Boston e Filadélfia, cinco relatórios por ano dos bancos nacionais e uma variedade de relatórios oficiais e privados de bancos organizados sob leis estaduais.

Alguma ideia sobre a quantidade de investimento e especulação em curso pode ser obtida a partir das transações da Bolsa de Valores de Nova York, do número de alvarás de construção concedidos, da quilometragem de ferrovias em construção e de itens de notícias menos sistemáticos sobre novos contratos firmados, emissões de títulos oferecidas ao público e afins.

Por último e mais importante, estão as perspectivas de lucros para as ferrovias, cujas receitas brutas e líquidas são publicadas e comentadas com interesse infalível. Os lucros da United States Steel Corporation ocupam provavelmente o segundo lugar na estima geral. Segue-se uma massa diversificada de informações fornecidas por relatórios que as grandes corporações envolvidas em vários ramos de mineração, manufatura e bancos optam por tornar públicos. O outro lado da moeda é refletido pelas estatísticas de falências compiladas semanalmente por duas grandes agências mercantis.

Embora longe de completa, esta lista de materiais para avaliar a tendência das condições de negócios é longa. De fato, é demasiado longa para o empresário médio. Compilar e analisar os dados disponíveis exige mais tempo, mais esforço, mais habilidade estatística ou mais capacidade analítica do que a maioria dos homens tem para gastar na tarefa. Assim, o homem de negócios típico ignora a evidência desconcertante e lê apenas as conclusões resumidas elaboradas pelo editor do seu jornal financeiro ou pela agência de previsão que ele assina. O fato de o estudo de barômetros de negócios e a previsão do tempo comercial se ter tornado ele próprio um negócio lucrativo oferece prova convincente, simultaneamente, da necessidade e da dificuldade de utilizar eficazmente materiais publicados para todos. É de tais especialistas, e não do comerciante e fabricante médio, que podemos esperar a melhoria e a disseminação da informação necessária como base para o aperfeiçoamento do controle social sobre o funcionamento da economia monetária.

Ora, estes previsores profissionais ligados às redações de jornais financeiros, casas de investimento e afins não acham os dados já à mão demasiado elaborados. Eles têm o tempo e a paciência; eles adquirirão o que agora lhes falta de técnica estatística e habilidade analítica para extrair a essência de grandes massas de dados. O que eles mais precisam para melhorar as suas previsões é de materiais mais extensos e mais confiáveis para trabalhar. Mas é também perfeitamente possível melhorar o uso que fazem dos dados já disponíveis.

Na secção seguinte, são apresentadas certas sugestões para construir novos barômetros de negócios e melhorar os antigos. Como elas resultaram de um esforço para compreender os ciclos recentes, estas sugestões podem ser úteis aos previsores profissionais de negócios, se não ao público em geral.

3. SUGESTÕES PARA MELHORAR OS BARÔMETROS DE NEGÓCIOS

A. Novos Barômetros Necessários

Entre as adições mais necessárias à lista de barômetros de negócios americanos estão as seguintes:

Um número de índice geral do volume físico do comércio poderia ser feito a partir de dados sobre a produção de certos produtos básicos, os embarques ou recebimentos de outros, os registos do comércio externo e fontes semelhantes. No Apêndice do seu "Purchasing Power of Money", Irving Fisher mostrou que muito material para este fim já é fornecido incidentalmente em documentos oficiais, e sem dúvida muito mais poderia ser obtido para uma investigação que não retrocedesse para além de 1900 ou 1905. Médias separadas devem ser estabelecidas para os grandes departamentos da indústria, uma vez que as diferenças entre a atividade relativa em diferentes linhas seriam frequentemente não menos significativas do que as mudanças computadas no total. Na medida do possível, estas subdivisões do número de índice do volume físico do comércio deveriam corresponder às do índice de preços por atacado.

O plano do Sr. Hull para obter relatórios sobre contratos firmados para obras de construção e a percentagem de trabalho executado sobre contratos antigos merece consideração cuidadosa. Uma vez que provém de um homem intimamente familiarizado com o negócio, a viabilidade do plano não pode ser levianamente negada. Poucos conjuntos de números dariam mais percepção sobre as condições de negócios quando a prosperidade estivesse prestes a entrar em crise ou quando a depressão estivesse a pavimentar o caminho para a prosperidade.

Um número de índice dos preços relativos das obrigações, e números correspondentes indicando mudanças nas taxas de juro sobre empréstimos de longo prazo, não seriam difíceis de preparar. Mesmo que isoladas, estas duas séries possuiriam grande valor como reflexo da atitude dos investidores; mas seriam ainda mais úteis se acompanhadas de dados sobre o montante de obrigações e notas de curto prazo colocadas no mercado por empresas comerciais e por governos, quer centrais quer locais.

Certos estados — nomeadamente Nova York e Massachusetts — começaram a compilar estatísticas de desemprego numa escala modesta. Mas não temos dados abrangentes deste tipo comparáveis aos da Grã-Bretanha, França e Alemanha. O seu valor, como índice não só do bem-estar entre os assalariados, mas também das mudanças na atividade dentro de indústrias importantes e mudanças na procura de bens de consumo, é tal que torna esta deficiência das estatísticas americanas motivo de grave preocupação.

O que mais se deseja de tudo são estatísticas que retratem as flutuações relativas dos custos e lucros. Infelizmente, as dificuldades teóricas e práticas na obtenção de tais números são especialmente grandes. Mas certamente cada extensão da autoridade pública sobre a atividade corporativa deve ser utilizada para exigir métodos uniformes de contabilidade como os que foram impostos às ferrovias interestaduais, e os relatórios obtidos pelo governo devem ser disponibilizados em alguma forma significativa para a informação do público de negócios.

Finalmente, a informação praticamente disponível na previsão das condições de negócios poderia ser materialmente aumentada pela publicação mais rápida de muitos dos dados fornecidos pelo governo. Em numerosos casos, números extremamente valiosos não são entregues à imprensa até que se tenham tornado questões de interesse histórico em vez de notícias atuais.

B. Melhoria dos Antigos Barômetros

Os números de índices dos preços das mercadorias no atacado seriam mais úteis se fossem computadas séries separadas para matérias-primas e para os artigos fabricados a partir delas, e se as matérias-primas fossem subdivididas em produtos agrícolas, animais, florestais e minerais. Diferenças nas flutuações destes diversos grupos seriam de grande assistência na determinação das causas, e portanto do significado, das mudanças no total geral. Além disso, um número de índice de mercadorias idênticas nos Estados Unidos, Grã-Bretanha, França e Alemanha facilitaria o esforço de seguir os cursos concomitantes dos ciclos econômicos em diferentes países e antecipar a reação das condições estrangeiras sobre as domésticas.

Os preços das ações devem ser computados segundo o plano do número de índice, em vez de na forma atual de preços médios reais das ações numa lista selecionada de corporações ou nos valores de mercado agregados de certas emissões de títulos. Para facilitar as comparações, a base escolhida para o número de índice das ações deve concordar com a escolhida para os preços das mercadorias. As ações distintamente de investimento devem ser separadas das favoritas especulativas, e as médias devem ser estabelecidas para as ferrovias, indústrias e serviços públicos. Através de uma seleção adequada de dados, as flutuações nos preços relativos médios das ações industriais poderiam ser feitas para refletir a sorte das empresas especialmente preocupadas com o fornecimento de equipamento industrial. Em contraste, as ações de serviços públicos seriam afetadas por uma procura de consumo relativamente estável, e as ações ferroviárias pela atividade dos negócios em geral.

Os relatórios semanais e mensais de compensações também seriam mais úteis se fossem acompanhados de números de índices que refletissem a magnitude relativa das mudanças nos montantes reais. Médias separadas destes números relativos deveriam ser fornecidas para os centros em que as operações financeiras, a atividade industrial e as condições agrícolas são os fatores dominantes.

As estatísticas bancárias seriam mais instrutivas se os montantes reais dos itens principais fossem suplementados por rácios não só de reservas em relação aos depósitos, mas também de empréstimos em relação aos depósitos. Outro conjunto de números relativos deveria ser feito para refletir a amplitude das flutuações nos itens principais. Se as associações de câmaras de compensação exigissem dos seus membros e tornassem públicos os retornos precisos dos recebimentos e embarques de moeda, a atividade dos negócios em várias seções poderia ser acompanhada com maior certeza. Finalmente, um dos pontos mais obscuros das condições de negócios atuais na América poderia ser esclarecido se as taxas de desconto sobre papel comercial de primeira classe nestes vários centros pudessem ser regularmente apuradas. Mas talvez não possam ser compiladas estatísticas de taxas de desconto perfeitamente comparáveis até que as qualidades variáveis do papel comercial agora em circulação tenham sido padronizadas, permitindo que os bancos nacionais aceitem letras à moda europeia.

Estender esta lista de sugestões para melhorar os números do tipo já publicado seria fácil; mas já foi dito o suficiente para deixar claro o caráter das mudanças desejáveis. Em geral, a necessidade é de uma discriminação mais cuidadosa entre dados dessemelhantes, agora frequentemente agrupados num único total; a recolha de novos centros de dados já publicados para Nova York; métodos mais uniformes de compilação para garantir a comparabilidade do que pretende ser números semelhantes; e a computação de flutuações relativas sobre uma base comum. Em muitos, se não em todos os casos, um conjunto duplo de números relativos é altamente desejável — um conjunto referindo-se aos montantes reais médios numa década fixa, o outro conjunto fazendo comparações com o período correspondente do ano anterior.

C. Dificuldades

Quando a discussão sobre como melhorar o nosso controle sobre o funcionamento da economia monetária é assim reduzida a detalhes práticos, a complexidade da tarefa torna-se patente. Quase nenhuma das sugestões para melhorar ou estender os índices das condições de negócios deixa de evocar vários obstáculos que dificultam a obtenção de dados fidedignos — a relutância dos interesses privados em divulgar informações; a diversidade de práticas de negócios em vários ramos e secções do país; as mudanças contínuas na organização de negócios; modificações na importância relativa de diferentes matérias-primas e ainda mais nos tipos e qualidades de produtos manufaturados; os quebra-cabeças técnicos de classificação e média estatística, etc. Em vista destas dificuldades, a perspetiva de uma melhoria rápida nos dados para a previsão de negócios não é tão brilhante como se poderia desejar.

O vigor dos esforços para superar as dificuldades dependerá em grande parte das exigências dos empresários por um melhor serviço. Hoje, a única classe que demonstra o sentido mais claro da utilidade de um levantamento estatístico abrangente do presente dos negócios como base para a previsão do futuro dos negócios é o especulador de ações. Mas muitos homens que preferem chamar-se investidores, e um número crescente de corretores, banqueiros, comerciantes, fabricantes, empreiteiros e afins, estão a tornar-se consumidores ativos de tais relatórios. Uma vez que se pode contar com estas classes para assinar os jornais e as agências confidenciais que lhes prestam o serviço mais satisfatório, a previsão de negócios tornar-se-á sem dúvida uma profissão mais extensa e fará o progresso que for possível sob a iniciativa privada estimulada pela concorrência.

No que diz respeito a muitos ramos de negócios, no entanto, informações que sejam simultaneamente fidedignas e abrangentes não podem ser obtidas pela iniciativa privada. Se o governo estenderá o âmbito das suas atividades atuais neste campo será provavelmente determinado principalmente por grandes questões de política pública. Pois a maioria dos números compilados pelo governo são subprodutos de medidas tomadas com outros fins em vista do que aumentar o conhecimento sobre o funcionamento da economia monetária. Tal como as estatísticas ferroviárias são um resultado incidental da lei do comércio interestadual, e as estatísticas bancárias da Lei Bancária Nacional, etc., assim futuras adições à produção estatística do governo serão feitas ou não conforme o controle público sobre os assuntos de negócios for estendido ou restringido.

Certamente a comunidade como um todo tem um profundo interesse no desenvolvimento deste ramo do trabalho do governo. Os empresários que estudam jornais financeiros estão preocupados principalmente em obter lucros ou evitar perdas decorrentes dos altos e baixos dos mercados. A comunidade, pelo contrário, está interessada em reduzir as perturbações que estes altos e baixos do mercado causam no processo de produção e distribuição de bens úteis. Mas enquanto este processo de produção e distribuição de bens úteis estiver subordinado ao processo de ganhar dinheiro, o interesse da comunidade em estabilizar o ritmo da atividade econômica pode ser promovido dando a todos os empresários as mesmas e melhores oportunidades possíveis para conhecer o presente e prever o futuro. No entanto, uma percepção vívida do que poderia ser realizado nesta linha para o bem-estar geral não é comum. Esforços mais diretos para aplicar as agências governamentais à correção do que é considerado serem maus resultados do empreendimento comercial apelam mais fortemente à maioria dos eleitores. Daí que sejam aqueles que desejam ver a forma atual de organização econômica aperfeiçoada em vez de fundamentalmente alterada que estão mais preocupados em pressionar a exigência de melhores relatórios governamentais sobre as condições de negócios.

V. A "Superfície Monetária das Coisas" e o Que Se Passa Por Baixo

A teoria precedente dos ciclos econômicos está preocupada principalmente com as fases pecuniárias da atividade econômica. Os processos descritos relacionam-se com mudanças nos preços, investimentos, margens de lucro, capitalização de mercado das empresas, créditos, manutenção da solvência e afins — tudo relacionado com o ato de ganhar dinheiro, em vez de com a produção de bens ou com a satisfação de necessidades. Apenas dois fatores não pecuniários comandam muita atenção — mudanças no volume físico do comércio e na eficiência do trabalho — e mesmo estes dois são tratados com referência ao seu impacto sobre os lucros presentes e prospectivos.

A razão para permanecer assim na "superfície monetária das coisas" ao analisar os ciclos econômicos, em vez de tentar penetrar sob os motivos que atuam na conduta econômica, é a razão exposta na Parte II, Capítulo II. A atividade econômica moderna é imediatamente animada e guiada, não pela busca de satisfações, mas pela busca de lucros. Portanto, os ciclos econômicos são distintamente fenômenos de caráter pecuniário, por oposição a um caráter industrial. Mergulhar abaixo das considerações de negócios relativas ao lucro e à perda, lidar com "renda psíquica" e "custo psíquico", ou mesmo lidar com a produção e o consumo físico noutras dimensões que não a pecuniária, é distorcer o problema. Pois os processos realmente envolvidos em trazer prosperidade, crises e depressão são os processos executados pelos empresários no esforço de ganhar dinheiro. Os empresários recusam-se a complicar os seus problemas indo atrás do dólar para aquilo que o dólar representa, e quem quiser compreender o que eles estão a fazer deve tratar a sua ação tal como ela é.

Mas se as causas dos ciclos econômicos que são importantes traçar residem quase inteiramente dentro da ordem pecuniária, as consequências de momento são questões de bem-estar humano. As formas como a prosperidade dos negócios, a crise e a depressão reagem sobre o bem-estar físico e mental da comunidade são tão numerosas, no entanto, que é viável mencionar apenas algumas das mais significativas.

Primeiro, a provisão feita para satisfazer as necessidades da comunidade — o volume físico da produção atual — torna-se mais abundante quando o negócio é próspero e mais escasso nos estágios iniciais da depressão. Mas nem a expansão nem o encolhimento nesta oferta de bens úteis são tão grandes como a expansão e o encolhimento nos valores pecuniários correspondentes. E a alteração que ocorre na produção da indústria é distintamente maior no que diz respeito aos bens de produção do que aos bens de consumo estáveis. Isto é, a quantidade de alimentos, vestuário e afins produzidos pela comunidade para o seu próprio uso corrente é mais estável do que os barômetros de negócios sugerem.

Sobre a distribuição desta oferta atual de bens úteis, os ciclos econômicos exercem uma forte influência, uma vez que produzem mudanças generalizadas tanto nos rendimentos monetários como no poder de compra dos dólares. A situação precária da família do assalariado na economia monetária consiste em grande parte no encolhimento do emprego causado pela depressão. As privações físicas, as ansiedades e as humilhações impostas a esta classe pela incapacidade de encontrar trabalho são não só elas próprias um mal grave, mas também fontes prolíficas de novos males — intemperança, prostituição, ociosidade crónica, o abandono de famílias e o raquitismo das crianças. Os lucros encolhem sem dúvida em maior proporção do que os salários, e muitas famílias que retiram o seu rendimento desta fonte são forçadas a adotar economias dolorosas e a suportar muita ansiedade, embora raramente sofram privações tão extremas como as dos assalariados sem trabalho. Por outro lado, a classe relativamente pequena de pessoas cujos rendimentos permanecem realmente fixos durante a depressão lucra com a redução do custo de vida; por esta vantagem material, no entanto, pagam um preço elevado em incerteza e na participação simpática no sofrimento dos outros.

Os ciclos econômicos também afetam o bem-estar material ao influenciarem a seleção de líderes de negócios, a centralização do poder econômico e o progresso da técnica industrial.

A prosperidade estimula o empreendimento e incentiva os homens a estabelecerem-se por conta própria. Mas, ao tornar mais fácil a sobrevivência de gestores inaptos por um tempo, reduz um pouco a eficiência econômica da comunidade. Mesmo os empresários mais capazes, sob a pressão da pressa, relaxam um pouco as suas precauções contra o desperdício. Entretanto, os investidores tornam-se mais inclinados para empreendimentos imprudentes, e uma proporção crescente da energia da sociedade é desperdiçada em empreendimentos não lucrativos. A crise e a depressão, pelo contrário, servem pelo menos para eliminar os gestores menos competentes, para impor uma atenção vigilante aos detalhes a todos e para tornar os investidores cautelosos.

As vendas judiciais e as reorganizações a que a depressão dá origem oferecem a melhor oportunidade para o aumento de fortunas já grandes e para a ascensão de magnatas de negócios já poderosos. Desta forma, a depressão promove a centralização do controle no mundo dos negócios. Mas, pelo contrário, muitas vezes enfraquece ou destrói alianças frouxamente cimentadas ou pools para a regulação da concorrência. E a promoção de grandes combinações entre empresas anteriormente independentes é geralmente empreendida nos estágios intermediários da prosperidade, quando os investidores estão otimistamente inclinados e antes de os mercados monetário e de títulos se terem tornado restritos.

Para o progresso da técnica industrial, no sentido da aplicação prática de melhorias já inventadas, a fase mais favorável do ciclo econômico é a recuperação. A depressão obriga os homens a procurar qualquer método viável de redução de custos; mas oferece pouco incentivo para o gasto imediato de grandes somas em melhorias. É a época em que as alterações são planeadas; a recuperação é a época em que são executadas em maior escala. A prosperidade é menos favorável, não por falta de fundos, mas por falta de tempo e atenção.

Em geral, a prosperidade é uma época de atividade extenuante, recompensada pelo conforto material e animada por grandes esperanças. Os seus principais inconvenientes sociais são o desperdício incidental à pressa, a extravagância alimentada pela afluência e pelo otimismo, a obsessão da atenção pelos interesses de negócios e as ansiedades que obscurecem os seus dias finais. Uma crise intensifica estas ansiedades, particularmente para os empresários e investidores. O tumulto abranda na depressão; mas o abrandamento traz o desânimo para aqueles cujos medos se concretizaram e deixa outros com uma perspetiva sombria, na melhor das hipóteses. Para os trabalhadores, é a época de maior sofrimento — de excesso de trabalho quando empregados e de privação quando na rua. Para estas desvantagens, a sua repressão do desperdício, o estímulo de planos para melhorias técnicas e a imposição de cautela relativamente aos investimentos são apenas compensações parciais.

Por breve que seja, esta afirmação de como os ciclos econômicos reagem sobre o bem-estar social é suficiente para sugerir a dupla personalidade adquirida pelos cidadãos da economia monetária. Ganhar dinheiro para o indivíduo, prosperidade nos negócios para a nação, são fins artificiais de esforço impostos pelas instituições pecuniárias. Por baixo de um residem as atividades impulsivas do indivíduo — o seu labirinto de reações instintivas parcialmente sistematizadas em desejos conscientes, conhecimento definido e esforços intencionais. Por baixo do outro residem os ideais vagos e conflituantes de bem-estar social que os membros de cada geração refazem segundo as suas próprias imagens. Neste obscuro mundo interior residem os motivos e significados fundamentais da ação e dele emergem os padrões flutuantes pelos quais os homens julgam o que para eles vale a pena.

A economia monetária não suplantou, mas atrelou estas forças. Sobre a atividade humana e os ideais humanos imprimiu o seu próprio padrão. Como facilitou a divisão do trabalho, como deu um toque pecuniário ao desejo de distinção, como mudou a base do poder político e deu origem a novas classes sociais — estes resultados da economia monetária são amplamente reconhecidos. Como ensinou os homens a pensar nos termos da sua própria lógica formal, eficiente dentro de certos limites mas árida quando levada a extremos, foi parcialmente trabalhado por autores como Simmel, Sombart e Veblen. Como as suas exigências técnicas sujeitam a atividade econômica a alternâncias contínuas de expansão e contração, é o que este livro pretendeu descrever em detalhe.

Sujeitos como os homens estão ao domínio de conceitos e ideais pecuniários, podem ainda julgar o funcionamento da economia monetária por padrões mais íntimos e mais vitais. Tornar estes últimos padrões claros, mostrar de que formas definidas a procura de lucros os transgride e conceber métodos viáveis de remediar estes maus resultados, é uma grande parte da tarefa da reforma social. A teoria econômica não se revelará de grande utilidade neste trabalho a menos que compreenda as relações entre as instituições pecuniárias que o homem civilizado está a aperfeiçoar, a natureza humana que ele herda de antepassados selvagens e as novas forças que a ciência lhe empresta. Tratar o dinheiro como um símbolo vazio que "não faz diferença senão por conveniência" é um hábito superado em superficialidade apenas pelo hábito contra o qual protesta — o de tratar o ato de ganhar dinheiro como o objetivo final do esforço.

\end{document}